\section{Identificación del Problema}

Colombia ha enfrentado dificultades recurrentes en la ejecución de proyectos públicos de infraestructura, especialmente en términos de retrasos y sobrecostos, fenómenos ampliamente documentados en estudios recientes basados en datos de contratación \cite{gomez2025data,gomez2025analyzing}. A este panorama se suma la persistencia de dinámicas de corrupción en la contratación pública, que deterioran la eficiencia del gasto y favorecen resultados de bajo impacto social \cite{sierra2025contratacion,lyra2022fraud}. En este contexto, aparecen los llamados ``elefantes blancos'', entendidos como obras costosas, inconclusas o de baja utilidad social, que terminan generando más cargas que beneficios para los territorios.

Aunque se han fortalecido mecanismos institucionales de control, buena parte de la vigilancia sigue operando de manera reactiva, es decir, cuando los problemas ya son visibles o el daño ya ocurrió. Este comportamiento limita la capacidad de prevención de las entidades públicas y de control fiscal \cite{ospina2024percepcion,sierra2025contratacion}.

Paralelamente, la disponibilidad de datos abiertos de contratación pública ha crecido de forma significativa. En particular, los registros transaccionales de plataformas como SECOP II permiten analizar variables de planeación, selección, ejecución y liquidación contractual. La evidencia reciente muestra que variables como el tipo de proceso, el valor del contrato y el número de oferentes son útiles para anticipar desviaciones de tiempo y costo \cite{gomez2025data,gomez2025analyzing}. Además, la literatura internacional confirma que los modelos de aprendizaje automático permiten identificar señales de riesgo en contratación pública, incluso con información limitada \cite{decarolis2022corruption,tas2024machine}.

Sin embargo, persiste un vacío aplicado en la identificación temprana de riesgo para contratos de infraestructura financiados con recursos del Sistema General de Regalías (SGR), integrando de manera sistemática variables contractuales y financieras desde etapas iniciales del proceso. Por ello, se plantea desarrollar una herramienta predictiva que permita priorizar contratos con alta probabilidad de presentar irregularidades o resultados críticos de ejecución, contribuyendo a migrar de un control reactivo a uno preventivo. Este enfoque fortalecería la toma de decisiones y promovería una gestión más eficiente, transparente y responsable de los recursos públicos.